\zhbox{54}{64}{235}{452}{10}{发生过载（OL）故障时，MP2984 通过 $I_{LIMIT}$ 环路调节来限制 $I_{OUT}$。若禁用 $I_{LIMIT}$，则逐周期限制 $I_{OUT}$。逐周期限流期间，若 IC 工作于升压模式，限制 SWC 峰值电流（$I_{PEAK}$）；若工作于降压模式，SWB 保持导通，直到 $I_L$ 达到 $I_{LIMIT\_VALLEY}$，然后开始下一周期。因此，所有工作模式下的 $I_L$ 均可控制。\par
\textbf{输出电压调节}\par
MP2984 通过 FB 引脚调节 $V_{OUT}$，将 $V_{FB}$ 与内部 $V_{REF}$ 比较。$V_{REF}$ 根据 VREF 寄存器位设置在300 mV至2.047 V之间。EA 的 COMP 输出控制流向输出的 $I_L$。\par
\textbf{开关频率（$f_{SW}$）和扩频（FSS）}\par
MP2984 通过2位 FSW 寄存器将 $f_{SW}$ 配置为200、300、400或600 kHz。大多数应用建议采用400 kHz。\par
MP2984 具有扩频（FSS）和频率抖动功能。将 Dither 位（0x02h，D[4]）设为1以使能扩频和频率抖动，设为0则禁用。扩频可降低特定频率处的峰值发射。\par
器件使用2 kHz三角波调制内部振荡器，扩频范围为 $\pm6\%$。}
\zhbox{211}{535}{42}{20}{7}{基准频率}
\zhbox{54}{624}{234}{16}{10}{\centering\textbf{图7：频率扩频}}
\zhbox{54}{642}{234}{34}{10}{扩频可配合200、300、400或600 kHz开关频率工作。}
\zhbox{324}{64}{235}{637}{10}{\textbf{栅极驱动与自举（BST）电源}\par
MP2984 为 H 桥 MOSFET 提供4个 N 沟道功率 MOSFET 栅极驱动器（见第24页图3）。每个驱动器均可输出及吸收电流。降压模式下 LG1 与 HG1 开关工作，HG2 保持导通；升压模式下 LG2 与 HG2 开关工作，HG1 保持导通。LG1、LG2 由 $V_{CC}$ 供电，HG1、HG2 分别由 $V_{BST1}$、$V_{BST2}$ 供电。\par
BST1 与 SW1 之间、BST2 与 SW2 之间必须连接供电电容，可通过内部二极管由 VCC 充电，也可彼此充电。\par
\textbf{过压保护（OVP）}\par
MP2984 监测反馈电压 $V_{FB}$。若 $V_{FB}>127\%V_{REF}$ 且 OVP\_MODE=01，IC 通过内部放电电阻对输出电容（$C_{OUT}$）放电；$V_{FB}$ 降至调节电压的111\%时停止放电。若 OVP\_MODE=00，即便 $V_{FB}$ 超过过压阈值，也没有用于停止开关工作的逻辑。若 OVP\_MODE=10，$V_{OUT}$ 达到 $127\%V_{REF}$ 时 IC 锁断。\par
\textbf{中断（INT）}\par
INT 为开漏输出，在过流、过压、过温或 $V_{OUT}$ 不在正常范围时指示中断。\par
若触发 OCP、输出 OVP 或热关断，对应寄存器位（OCP：D[1]，OVP：D[2]，OTP：D[4]）置1。INT 根据对应屏蔽寄存器设置拉低，以指示中断。器件关闭时 INT 为开漏状态。\par
\textbf{电流监测输出（IMON）}\par
MP2984 通过检测电阻检测平均负载电流，在 IMON 引脚输出电压信号。该信号为 IAVGP 与 IAVGN 间电压差（$V_{IAVGP}-V_{IAVGN}$）经放大后的信号。在 IMON 与 AGND 间接一个小电容。IMON 电压（$V_{IMON}$）可由式(4)计算：}
